\documentclass[11pt]{article}
\usepackage{graphicx}
\usepackage{amssymb}
\usepackage{bm}
\usepackage{mathtools}
\usepackage{amsmath}
\usepackage{commath}
\usepackage{amsthm}
\usepackage{enumitem}
\usepackage{float} 
\usepackage{array}
\usepackage{outlines}
\usepackage{setspace}
\usepackage[colorinlistoftodos]{todonotes}
\definecolor{green}{rgb}{0.4, 0.69, 0.2}
\definecolor{lavender(floral)}{rgb}{0.71, 0.49, 0.86}
\definecolor{asparagus}{rgb}{0.53, 0.66, 0.42}
\definecolor{blue}{rgb}{0.0, 0.53, 0.74}
\setlist[enumerate,1]{label=\roman*)}
\setlist[enumerate,2]{label=\alph*)}
\newtheorem{theorem}{Theorem}
\newtheorem{corollary}{Corollary}
\newtheorem{lemma}{Lemma}
\newtheorem{remark}{Remark}
\newtheorem{definition}{Definition}
\newtheorem{prop}{Proposition}
\newcommand*\diff{\mathop{}\!\mathrm{d}}
\usepackage{tikz}
\usetikzlibrary{arrows,shapes,trees, calc}


%%% *** PREAMBLE *** %%%
\title{Existence --- the strong solution of the downwind-matching Primitive Equations}

\date{}

%%%***  CONTENT *** %%%
\begin{document}

\maketitle

\begin{abstract}

Based on "Local and Global Well-Posedness of Strong Solutions to the 3D Primitive Equations with Vertical Eddy Diffusivity".

\end{abstract}

\newpage
\tableofcontents
\newpage

\section{The main steps}

We need an existence result for the system (E-PEs): extended Primitive Equations.

    \begin{equation} \label{HNSC1-dw}    
    \partial_t u_1 + \mathbf{u} \cdot \nabla u_1 - \Delta_\nu u_1 - \gamma u_2 + \partial_1 p = 0 \text{ in } \Omega \times (0,T)
    \end{equation}
    
    \begin{equation} \label{HNSC2-dw}   
    \partial_t u_2 + \mathbf{u} \cdot \nabla u_2 - \Delta_\nu u_2 + \gamma u_1 + \partial_2 p = 0 \text{ in } \Omega \times (0,T)
    \end{equation}
    
    \begin{equation} \label{HNSC3-dw}   
    \partial_3 p = 0 \text{ in } \Omega \times (0,T)
    \end{equation}
    
    \begin{equation} \label{HNSC4-dw}   
    \nabla \cdot \mathbf{u} = 0 \text{ in } \Omega \times (0,T)
    \end{equation}
    
    \begin{equation} \label{HNSC5-dw}   
    \partial_t c + \mathbf{u} \cdot \nabla c = K_2 \partial_{22} ^ 2 c + K_3 \partial_{33} ^2 c + s \text{ in } \Omega \times (0,T)
    \end{equation}
    
Structurally, this is "almost" identical to the Primitive Equations with $T$ as a temperature, except there we have the additional connection between the $T$ function and the pressure:

    \begin{equation} \label{HNSC1-T}    
    \partial_t u_1 + \mathbf{u} \cdot \nabla u_1 - \Delta_\nu u_1 - \gamma u_2 + \partial_1 p = 0 \text{ in } \Omega \times (0,T)
    \end{equation}
    
    \begin{equation} \label{HNSC2-T}   
    \partial_t u_2 + \mathbf{u} \cdot \nabla u_2 - \Delta_\nu u_2 + \gamma u_1 + \partial_2 p = 0 \text{ in } \Omega \times (0,T)
    \end{equation}
    
    \begin{equation} \label{HNSC3-T}   
    \partial_3 p + T = 0 \text{ in } \Omega \times (0,T)
    \end{equation}
    
    \begin{equation} \label{HNSC4-T}   
    \nabla \cdot \mathbf{u} = 0 \text{ in } \Omega \times (0,T)
    \end{equation}
    
    \begin{equation} \label{HNSC5-T}   
    \partial_t T + \mathbf{u} \cdot \nabla T =  K_3 \partial_{33} ^2 T + Q \text{ in } \Omega \times (0,T)
    \end{equation}
    
We need to examine what differences the change from (\ref{HNSC3-T}) to (\ref{HNSC3-dw}) brings over the course of the proof.

The main difference seems to be what shows after reformulation. In C-L-T's article, (1.23),

$$ \nabla_H \bigg (\int_{-h}^z T (x,z,\xi,t ) \diff \xi \bigg) $$

appears in the horizontal momentum equation. $\nabla_H T$ is exactly the part that otherwise the convection-diffusion equation does not have information about directly. We need to check if this causes a change.

It does in the application process of $\theta,$ where $u$ denotes $u = \partial_z v,$ and $ \theta = \nabla_H \cdot u + R_1 T.$ In their case this turns out to be a function definition that has full $H^2$ regularities in the sense of Proposition 5.3 and 5.4, because when the $L_1 = - \frac{1}{R_1} \Delta_H - \frac{1}{R_2} \partial_z^2$ full Laplacian and a time derivative are applied to the function $\theta$, then the horizontal temperature Laplacians from $\nabla_H \cdot u$ and $R_1 T$ cancel each other.
In more detail, $\partial_t \theta + L_1 \theta$ will contribute:

\begin{itemize}

  \item from $\partial_t \nabla_H \cdot u,$ we will have
  $$\nabla_H \cdot \nabla_H \partial_z \bigg (\int_{-h}^z T (x,z,\xi,t ) \diff \xi \bigg), $$
  
  \item from $L_1 R_1 T$ we will have
  $$ - R_1 \frac{1}{R_1} \Delta_H T,$$
  
\end{itemize}

which cancel out perfectly. This cancellation makes it possible that on p71- p72 they only arrive to have $\norm{\partial_z T}_{H^2}^2$ in the estimates, but not a full $H^3$ norm of $T$ ($\int_0^t \norm{\partial_z T}_{H^2}^2 \diff s$ is bounded, but we do not have an analogous result for the complete $H^3$ norm.)

Thus, in our case we can not keep $\theta$'s definition including the $T$ function (for their case it made sense to create $\theta$ as it would have the necessary regularities, and the momentum equation contained $T$) --- on the other hand, splitting $T$ off of $\theta$ also alters the proof significantly. The structure described on p58-p60 works in a simultaneous way for $T$ and $v.$ The general result $(4.10)$ is used later also for $T.$ If (4.10) does not include $T$ anymore, then for $T$ we need to create a solution on its own.

This suggests the following idea:

\begin{enumerate}

  \item Follow the local existence result without any particular alterations (except for leaving out the $T$ function from the third equation) and verify that it is still valid.
  
  \item As for the global existence, split $T$ from $\theta,$ and follow a sequential method instead of a simultaneous one:
  
  \begin{itemize}
    \item verify the global existence for the velocity part only (the available results on the existence of the strong solutions of the Primitive Equations are for the temperature-including version)
    \item using the already proven $v$ and its regularities, show a uniform estimate for the energy of $T.$
  \end{itemize}

\end{enumerate}

\section{Step 1. The Regularised System with Full Diffusion}

Let's consider

$$ \partial_t v + L_1 v + (v \cdot \nabla_H) v - \bigg ( \int_{-h}^z \nabla_H \cdot v \bigg ) \partial_z v+ f_0 k \times v + \nabla_H p_s(x,y,t) = 0$$

$$ \nabla_H \cdot \bar{v} = 0 $$

$$ \partial_t C + L_2 C - \epsilon \Delta_H C + v \cdot \nabla_H C - \bigg ( \int_{-h}^z \nabla_H \cdot v \bigg ) \bigg ( \partial_z C + \frac{1}{h} \bigg).$$

The system is formally analogous to (2.1)-(2.3):

$$ \partial_t v + L_1 v + (v \cdot \nabla_H) v - \bigg ( \int_{-h}^z \nabla_H \cdot v \bigg ) \partial_z v+ f_0 k \times v + \nabla_H \bigg ( p_s(x,y,t) -  \int_{-h}^z T \diff \xi \bigg )= 0$$

$$ \nabla_H \cdot \bar{v} = 0 $$

$$ \partial_t T + L_2 T - \epsilon \Delta_H T + v \cdot \nabla_H T - \bigg ( \int_{-h}^z \nabla_H \cdot v \bigg ) \bigg ( \partial_z T + \frac{1}{h} \bigg),$$

except we do not have the concentration term in the momentum equation. In some way it is similar to taking $T_1$ and $T_2$ in (2.1)-(2.3) instead of taking one common $T.$ In (2.1) we use $T_1 = 0,$ in (2.3) we use $T_2=C.$

The $\mathcal{F}$ map remains the same, with the modification of the left-hand side functions $A(v,T)$ and $B(v,T).$ These become $A(v), B(v)$ respectively, by leaving out the $T-$ including terms from (2.16) and (2.17). The periodicity and parity remains, and their regularities remain.

Because of the unchanged regularities and properties, the right-hand side of (2.8)-(2.11) stays $L^2_t H^2_x$ regular and we can still use the standard $L^2$ theory \todo{what does it say exactly?} of linear Stokes equations and parabolic equations and have a $(U,V,\mathcal{T})$ solution. (check (2.19))

Analogously, $D(v,T)$ will become $D(T),$ keeping only the $v$-dependent terms. The properties connecting $A(v), B(v)$ and $D(v)$ still hold.

\textbf{Proposition 2.2}

Estimating the functions $D(v)$ and $E(v,T)$ we arrive eventually to the same estimate for $I(t)$ on p45. Note that the calculations for $E(v,T)$ are identical, while for $D(v)$ we can leave out some non-negative terms from the right-hand side of the estimate, namely do not have the $\norm{T_1 - T_2}_{H^2}^2$ term in (2.20) --- however this does not lead to any change in $I(t),$ since $E(v,T)$ does contain $C \norm{T_1 - T_2}_{H^2}^2$.

The rest of the proof holds the same, keeping in mind that we have $D(v_i)$ where they use $D(v_i, T_i).$

\textbf{Proposition 2.3} and \textbf{Proposition 2.1} are general in nature and hold without modification.

\section{Step 2. Local existence and uniqueness}

\section{Step 3. Global existence}

\section{Questions}

\begin{itemize}

  \item As for the odd / even structure. It seems that this is used because for passing from $(-h, 0)$ to $(-h, h)$ we need to extend the function in a way to a domain where it did not exist before. In order to define it somehow in a way that afterwards the restriction will work as well when going backwards to $(-h, 0)$ it simply is a natural way to extend. It is natural because the odd and even properties are invariant for the system, and using symmetries we can extend the function in a way that afterwards the restriction will work for the original one. The question is, what is the benefit from passing from $(-h, 0)$ to $(-h, h).$

  \item If the (\ref{HNSC1-dw}) - (\ref{HNSC5-dw}) equations do have a unique strong solution $(v, C)$ with $(v_0, C_0)$ then taking this $v$, and combining it with the same $v_0$, and $Q = 0, T_0 = 0, T \equiv 0,$ then the resulting $(v,T) = (v,0)$ will be a solution for (\ref{HNSC1-T}) - (\ref{HNSC5-T}) with $(v_0, T_0),$ simply because by definition they satisfy the equations. Then, using the uniqueness of the solution of (\ref{HNSC1-T}) - (\ref{HNSC5-T}), it means that the only solution for this system with $T_0 = 0, Q = 0$ is $T \equiv 0.$ This is a bit surprising, because even though the source and the initial data is zero, through the velocity fluctuations the pressure changes, and with that the pressure should change too, shouldn't it?
  
  \item After all, why $p$ isn't a part of a solution as $T$ and $v$? The $T$-free hydrostatic scenario $p$ fell out due to the divergence-free condition in many places, but here it is even present in the reformulated version.

  \item We have
  
    $$\partial_z w + \nabla_H \cdot v = 0,$$
    
    $$ \partial_z p + T = 0.$$
  
  These lead to
  
    $$ w = - \int_{-h}^z \nabla_H \cdot v, $$
    
    $$ p = p_s(x,y,t) - \int_{-h}^z T .$$

Why the difference? They are structurally the same, in terms of containing information about a function through its $\partial_z.$ Why does $p$ have a $p_s(x,y)$ part, and $w$ does not?

\end{itemize}

\end{document}

